% Generated draft: requires mathematical and source review.
\documentclass[11pt]{article}
\usepackage[margin=25mm]{geometry}
\usepackage{fontspec}
\setmainfont{Latin Modern Roman}
\usepackage{amsmath,amssymb,mathtools,unicode-math}
\setmathfont{Latin Modern Math}
\usepackage{enumitem,needspace,xurl,hyperref}
\setlength{\parindent}{0pt}
\setlength{\parskip}{5pt}
\setlength{\emergencystretch}{3em}
\newcommand{\N}{\mathbb{N}}
\newcommand{\Z}{\mathbb{Z}}
\newcommand{\Q}{\mathbb{Q}}
\newcommand{\R}{\mathbb{R}}
\newcommand{\C}{\mathbb{C}}
\newcommand{\F}{\mathbb{F}}
\newcommand{\A}{\mathbb{A}}
\newcommand{\PP}{\mathbb{P}}
\newcommand{\E}{\mathbb{E}}
\newcommand{\eps}{\varepsilon}
\newcommand{\dd}{\,\mathrm{d}}
\begin{document}
\section*{UnsolvedMath ID: 30001857}
\textbf{Source problem code:} catalogue entry
\subsection{Definitions and mathematical statement}
Let $\F_2=\Z/2\Z$. A binary linear code of length $72$ is a linear subspace $C\subseteq\F_2^{72}$. For $c=(c_1,\ldots,c_{72})$, define its Hamming weight by $\operatorname{wt}(c)=|\{i:c_i\ne0\}|$. Define
\[
C^\perp=\left\{x\in\F_2^{72}:\sum_{i=1}^{72}x_ic_i=0\text{ in }\F_2\text{ for every }c\in C\right\}.
\]
The code is self-dual if $C=C^\perp$, and doubly even if $\operatorname{wt}(c)\equiv0\pmod4$ for every $c\in C$. Its minimum distance is
\[
d(C)=\min\{\operatorname{wt}(c):c\in C\setminus\{0\}\}.
\]
The parameters $[72,36,16]$ specify length $72$, dimension $36$ over $\F_2$, and minimum distance exactly $16$; this is the extremal case intended by the supplied record.

\textbf{Existence question.} Does there exist a linear subspace $C\subseteq\F_2^{72}$ such that
\[
\dim_{\F_2}C=36,\qquad C=C^\perp,\qquad
\operatorname{wt}(c)\equiv0\pmod4\quad(\forall c\in C),\qquad d(C)=16?
\]
The original abbreviated question asks only for an extremal code of length $72$. The binary, self-dual, doubly even hypotheses and the parameters above follow the supplied corrected statement, whose claimed source verification has not been independently checked here.
\subsection{Short English statement}
Does there exist a $36$-dimensional binary linear code of length $72$ that equals its orthogonal complement, has every codeword weight divisible by $4$, and has minimum nonzero weight exactly $16$?
\subsection{Sources}
\begin{itemize}
\item Local UnsolvedMath registry, \texttt{data/unsolvedmath/problems.json}, permanent ID $30001857$, source problem code \texttt{catalogue entry}, in the Oberwolfach Reports Open Problems collection. The record gives proposed year $2011$ and names no individual proposer.
\item Supplied report citation: \emph{Computational Group Theory (2012)}, Oberwolfach Reports, pp.~2113--2161, DOI \texttt{10.4171/owr/2011/37}. \url{https://doi.org/10.4171/owr/2011/37}. The supplied background describes Oberwolfach Reports as documenting mathematical workshops and their open-problem discussions. The report-title year $2012$ and the supplied proposed year $2011$ are retained without resolving their bibliographic relationship.
\item Supplied literature source labeled ``Borello'': \url{https://arxiv.org/abs/1203.3321}.
\item Supplied literature source labeled ``OEIS open problem record'': \url{https://oeis.org/A322339}.
\item Supplied literature source labeled ``OWR report 37'': \url{https://doi.org/10.4171/owr/2011/37}.
\item Literature assessment dated \texttt{2026-08-21}: the supplied curation labels the existence question open and states that an extremal binary self-dual doubly even $[72,36,16]$ code remains unknown. This is supplied dated curation derived from the OWR record, not an independently verified literature assessment. Time-sensitive status and cited claims require independent verification.
\end{itemize}
\end{document}
